% Experiment 11: Implementation and Analysis of Feistel Cipher Structure
\section{Theory}
The Feistel cipher structure is a symmetric block cipher design used as the
foundation of many well-known encryption algorithms, including DES. It
divides the input block into two equal halves, $L_0$ and $R_0$, and applies
a sequence of identical rounds. In each round $i$, the right half is
processed by a round function $F$ together with a round-specific subkey
$K_i$, and the result is XORed with the left half:
\[
L_i = R_{i-1}, \qquad R_i = L_{i-1} \oplus F(R_{i-1}, K_i)
\]
A key advantage is that $F$ need not be invertible: decryption simply reruns the same structure with subkeys in reverse order, allowing arbitrarily complex non-invertible functions while keeping the overall cipher reversible.

\section{Algorithm}
\begin{enumerate}
  \item Split the plaintext block into two equal halves, $L_0$ and $R_0$.
  \item Generate a set of round keys $K_1, K_2, \dots, K_n$ from the master key.
  \item For each round $i = 1$ to $n$:
  \begin{enumerate}
    \item Compute $F(R_{i-1}, K_i)$.
    \item Set $L_i = R_{i-1}$.
    \item Set $R_i = L_{i-1} \oplus F(R_{i-1}, K_i)$.
  \end{enumerate}
  \item After the final round, concatenate $L_n$ and $R_n$ to form the ciphertext.
  \item To decrypt, apply the same round structure to the ciphertext, using
        the round keys in reverse order ($K_n, K_{n-1}, \dots, K_1$).
  \item Verify that the decrypted output matches the original plaintext.
\end{enumerate}

\section{Code}
\begin{minted}{cpp}
#include<bits/stdc++.h>
using namespace std;
class feistel {
  private:
    int rounds;
    vector<int> roundKeys;
    static const int HALF_BITS = 4;
    int roundFunction(int R, int K) {
      int x = (R ^ K) & ((1 << HALF_BITS) - 1);
      return ((x << 1) | (x >> (HALF_BITS - 1))) & ((1 << HALF_BITS) - 1);
    }
  public:
    feistel(int masterKey, int r);
    int encrypt(int plaintext);
    int decrypt(int ciphertext);
};
feistel::feistel(int masterKey, int r) {
  rounds = r;
  masterKey &= (1 << HALF_BITS) - 1;
  for (int i = 0; i < rounds; i++) {
    roundKeys.push_back((masterKey + i * 3) & ((1 << HALF_BITS) - 1));
  }
  cout << "Round keys generated: ";
  for (int k : roundKeys) cout << bitset<HALF_BITS>(k) << " ";
  cout << endl;
}
int feistel::encrypt(int plaintext) {
  int L = (plaintext >> HALF_BITS) & ((1 << HALF_BITS) - 1);
  int R = plaintext & ((1 << HALF_BITS) - 1);
  for (int i = 0; i < rounds; i++) {
    int newL = R;
    int newR = L ^ roundFunction(R, roundKeys[i]);
    L = newL;
    R = newR;
    cout << "  Round " << i + 1 << ": L=" << bitset<HALF_BITS>(L)
         << " R=" << bitset<HALF_BITS>(R) << endl;
  }
  return (L << HALF_BITS) | R;
}
int feistel::decrypt(int ciphertext) {
  int L = (ciphertext >> HALF_BITS) & ((1 << HALF_BITS) - 1);
  int R = ciphertext & ((1 << HALF_BITS) - 1);
  for (int i = rounds - 1; i >= 0; i--) {
    int newR = L;
    int newL = R ^ roundFunction(L, roundKeys[i]);
    L = newL;
    R = newR;
    cout << "  Round " << (rounds - i) << ": L=" << bitset<HALF_BITS>(L)
         << " R=" << bitset<HALF_BITS>(R) << endl;
  }
  return (L << HALF_BITS) | R;
}
int main() {
  cout << "Enter plaintext (0-255): ";
  int plaintext;
  cin >> plaintext;
  cout << "Enter master key (0-15): ";
  int masterKey;
  cin >> masterKey;
  plaintext &= 255;

  feistel cipher(masterKey, 4);

  cout << "\n--- Encryption ---" << endl;
  cout << "Plaintext : " << bitset<8>(plaintext) << endl;
  int cipherText = cipher.encrypt(plaintext);
  cout << "Ciphertext: " << bitset<8>(cipherText) << endl;

  cout << "\n--- Decryption ---" << endl;
  int decrypted = cipher.decrypt(cipherText);
  cout << "Decrypted : " << bitset<8>(decrypted) << endl;

  return 0;
}
\end{minted}

\section{Result and Discussion - I/O}
We take integer values as input as they are easiest to reference as bit stream and are easy to run xor and other mathematical operations. 
\begin{verbatim}
Enter plaintext (0-255): 170
Enter master key (0-15): 5
Round keys generated: 0101 1000 1011 1110

--- Encryption ---
Plaintext : 10101010
Round 1: L=1010 R=0101
Round 2: L=0101 R=0001
Round 3: L=0001 R=0000
Round 4: L=0000 R=1100
Ciphertext: 00001100

--- Decryption ---
Round 1: L=0001 R=0000
Round 2: L=0101 R=0001
Round 3: L=1010 R=0101
Round 4: L=1010 R=1010
Decrypted : 10101010
\end{verbatim}


\section{Conclusion}
We successfully implemented and tested an 8-bit Feistel cipher with an XOR-and-rotate round function. The experiment confirms that Feistel networks allow complex round functions without requiring their inversion during decryption.
